\section*{第 \ref{ch:b2:acyl-substitution} 章——羧酸衍生物}

\begin{solution}{exo:b2:acyl-substitution:1}
乙酰胺 $<$ 乙酸乙酯 $<$ 乙\omterm{def:b2:acyl-substitution:derivative}{酸酐} $<$ 乙\omterm{def:b2:acyl-substitution:derivative}{酰氯}。
\end{solution}

\begin{solution}{exo:b2:acyl-substitution:2}
乙酸（+ \ce{HCl}）；乙酸乙酯；乙酰胺（+ 氯化铵，第二当量的氨用于捕获
\ce{HCl}）；乙\omterm{def:b2:acyl-substitution:derivative}{酸酐}（+ \ce{NaCl}）；
\textit{N},\textit{N}-二甲基乙酰胺（+ 氯化三乙基铵）。
\end{solution}

\begin{solution}{exo:b2:acyl-substitution:3}
\ce{C6H5COOCH3 + NaOH -> C6H5COONa + CH3OH}。$13.6/136.15 = \qty{0.0999}{mol}$，
需要等量的 \ce{NaOH}：\qty{4.00}{g}。
\end{solution}

\begin{solution}{exo:b2:acyl-substitution:4}
$\gamma$-丁\omterm{def:b2:acyl-substitution:lactone}{内酯}（氧杂环戊-2-酮）：由四个碳和一个氧构成的五元环。
\end{solution}

\begin{solution}{exo:b2:acyl-substitution:5}
氨加成到羰基碳上；质子由氮转移到氧上（或转移到
离去基团上）；乙醇盐（质子转移后以乙醇形式）离去，\ce{C=O} 重新形成：
生成乙酰胺和乙醇。
\end{solution}

\begin{solution}{exo:b2:acyl-substitution:6}
酚羟基 \ce{OH} 被酰化（酸催化）：2-乙酰氧基苯甲酸；副产物为乙酸。
\end{solution}

\begin{solution}{exo:b2:acyl-substitution:7}
2-甲基丙-2-醇。第一次加成生成丙酮，其亲电性强于酯，
立即再接受第二个甲基。
\end{solution}

\begin{solution}{exo:b2:acyl-substitution:8}
丁-1,4-二醇：酯被还原为两个伯醇，环随之被打开。
\end{solution}

\begin{solution}{exo:b2:acyl-substitution:9}
一当量时：$x^2/(1 - x)^2 = 4$，$x = 2/3$。三当量时：$x^2/[(1 - x)(3 - x)] = 4$，$3x^2 - 16x
+ 12 = 0$，$x = 0.90$。随生成随移除水（迪安–斯塔克装置）可使反应进一步进行。
\end{solution}

\begin{solution}{exo:b2:acyl-substitution:10}
\ce{SOCl2} 生成苯甲\omterm{def:b2:acyl-substitution:derivative}{酰氯}；与两当量二乙胺（或一当量二乙胺加三乙胺）
反应生成酰胺。酸与胺直接混合则发生酸碱反应，生成苯甲酸
二乙基铵，只在强烈加热时才得到酰胺。
\end{solution}

\begin{solution}{exo:b2:acyl-substitution:11}
三油酸甘油酯每分子需三个 KOH，故每克需：$3 \times 56.11/885.4 = \qty{0.190}{g}$，\omterm{def:b2:acyl-substitution:saponification}{皂化}值为
190.
\end{solution}

\begin{solution}{exo:b2:acyl-substitution:12}
\ce{NaBH4}：只还原酮，生成仲醇。\ce{LiAlH4}：将酮还原为
仲醇，将酯还原为伯醇（并释放出醇部分），将酰胺还原为胺。
\end{solution}

\begin{solution}{pb:b2:acyl-substitution:1}
\textbf{1.} 油酸 \ce{C18H34O2}；甘油 \ce{C3H8O3}。
\textbf{2.} \ce{C57H104O6}: $57 \times 12.011 + 104 \times 1.008 + 6 \times 15.999 = \qty{885.45}{g/mol}$.
\textbf{3.} 三个。
\textbf{4.} 顺式双键使链弯折，堆积不紧密：分子间作用力弱，
熔点较低。
\textbf{5.} 甘油和三分子油酸钠，即一种肥皂。
\textbf{6.} 脂肪酸甲酯的混合物。
\textbf{7.} \ce{C57H104O6 + 3CH3OH -> C3H8O3 + 3C19H36O2}.
\textbf{8.} 甲醇盐离子由甲醇与碱生成，其亲核性远强于
甲醇；它进攻酯的羰基。
\textbf{9.} 一个酯基的羰基碳（带 \ce{O-}）、甲醇盐的 \ce{OCH3} 以及
甘油基的氧：四面体形。
\textbf{10.} \omterm{def:b2:acyl-substitution:transesterification}{酯交换}是一个平衡常数接近 1 的平衡；过量甲醇使其
向甲酯方向移动。
\textbf{11.} 甘油以下层形式分离：移去生成物可使平衡右移。
\textbf{12.} 否：每一步交换中甲醇盐均会再生。
\textbf{13.} 它中和了碱，生成油酸钠（肥皂）：催化剂被消耗。
\textbf{14.} 它使酯\omterm{def:b2:acyl-substitution:saponification}{皂化}：氢氧根（由水和甲醇盐生成）将酯切为肥皂和
甲醇，且不可逆。
\textbf{15.} 它消耗碱，使混合物乳化，并阻止甘油层分离。
\textbf{16.} 先以甲醇进行酸催化\omterm{def:b2:acyl-substitution:esterification}{酯化}，使游离酸
转化为甲酯。
\textbf{17.} 水使酯水解，并将催化剂转化为氢氧根，进而发生\omterm{def:b2:acyl-substitution:saponification}{皂化}。
\textbf{18.} $10^6/885.45 = \qty{1129}{mol}$.
\textbf{19.} $3 \times 1129 \times 32.04 = \qty{108.6}{kg}$.
\textbf{20.} $3 \times 1129 \times 296.50 = \qty{1004.6}{kg}$.
\textbf{21.} 输入：$1000 + 108.6 = \qty{1108.6}{kg}$；输出：$104.0 + 1004.6 = \qty{1108.6}{kg}$。
\textbf{22.} 油酸 \qty{20}{kg}，$20\,000/282.47 = \qty{70.8}{mol}$，需要等量的 \ce{NaOH}：
\qty{2.83}{kg}.
\textbf{23.} 每吨油脂可得甘油：$1129 \times 92.09 = \boldsymbol{\approx \qty{104}{kg}}$。
\end{solution}
